% GENERATED FILE. Edit canonical lessons, then run npm run generate:books.
% canonical-source-hash: fnv1a64:7c907c57d3b3d98c
% canonical-lessons: SA-C385-spardhate, SA-C385-pratidvandvi, SA-C385-jeta, SA-C385-puraskaroti, SA-C385-utkarsah

\chapter{Competing}
\label{ch:sa-competing}
\hlchaptermodality{385}

\begin{chapteropening}
\textbf{By the end of this chapter:} \emph{I can say who will win if they compete well.}

Complete the last lesson of chapter 385: I can say who will win if they compete well.
\end{chapteropening}

\section[\texorpdfstring{\sk{स्पर्धते} (spardhate)}{स्पर्धते (spardhate)}]{\sk{स्पर्धते} (spardhate) — competes}

\label{lesson:SA-C385-spardhate}



\begin{quote}
Before the new one: say the Sanskrit for carries out, then the Sanskrit for encourages.
\end{quote}

\subsection*{You'll want to know: \sk{स्पर्धते}}
\textbf{\sk{स्पर्धते}} — \emph{spardhate} — \textquotedblleft{}competes\textquotedblright{}.

\begin{quote}
\sk{यदि} \sk{मीरा} \sk{स्पर्धते}, \sk{तर्हि} \sk{सा} \sk{प्रथमा} \sk{भविष्यति।}
\end{quote}

\begin{quote}
\emph{yadi mīrā spardhate, tarhi sā prathamā bhaviṣyati.}
\end{quote}

\begin{quote}
If Mira competes, she will come first.
\end{quote}

Past (she …): \sk{सा} \sk{स्पर्धितवती}, \emph{sā spardhitavatī}.

\begin{cousinweb}
\textbf{\sk{स्पर्धते}} (\emph{spardhate}) comes from the root √spardh, \textquotedblleft{}to contend, to vie\textquotedblright{}; MW gives \emph{spardhita} \textquotedblleft{}contending or competing together\textquotedblright{} (MW).
\end{cousinweb}

\begin{grammarlens}[title={What you've built}]
One more word for saying what happens if.
\end{grammarlens}

\subsection*{Your turn}
\begin{itemize}
\raggedright
  \item \emph{Say it:} \emph{spardhate}
  \item \emph{Say it:} \emph{spardhate}, once more
  \item \emph{Say it:} \emph{yadi mīrā spardhate, tarhi sā prathamā bhaviṣyati.} — If Mira competes, she will come first.
  \item \emph{Recall:} say the Sanskrit for carries out, then the Sanskrit for encourages, then say \emph{spardhate} again
\end{itemize}

\begin{tcolorbox}[breakable,colback=teal!4,colframe=teal!35!black,title={Before you move on}]
What does \sk{स्पर्धते} mean? (\textquotedblleft{}Competes\textquotedblright{}.) Say it once more.
\end{tcolorbox}
\section[\texorpdfstring{\sk{प्रतिद्वन्द्वी} (pratidvandvī)}{प्रतिद्वन्द्वी (pratidvandvī)}]{\sk{प्रतिद्वन्द्वी} (pratidvandvī) — a rival}

\label{lesson:SA-C385-pratidvandvi}



\begin{quote}
Before the new one: say the Sanskrit for encourages, then the Sanskrit for competes.
\end{quote}

\subsection*{You'll want to know: \sk{प्रतिद्वन्द्वी}}
\textbf{\sk{प्रतिद्वन्द्वी}} — \emph{pratidvandvī} — \textquotedblleft{}a rival\textquotedblright{}.

\begin{quote}
\sk{यदि} \sk{प्रतिद्वन्द्वी} \sk{बलवान्}, \sk{तर्हि} \sk{अधिकः} \sk{अभ्यासः} \sk{आवश्यकः।}
\end{quote}

\begin{quote}
\emph{yadi pratidvandvī balavān, tarhi adhikaḥ abhyāsaḥ āvaśyakaḥ.}
\end{quote}

\begin{quote}
If the rival is strong, more practice is needed.
\end{quote}

A masculine noun.

\begin{cousinweb}
\textbf{\sk{प्रतिद्वन्द्वी}} (\emph{pratidvandvī}) is built on \emph{prati}, \textquotedblleft{}against\textquotedblright{}, and \emph{dvandva}, \textquotedblleft{}a pair, a strife\textquotedblright{}; MW gives \emph{prati-dvaṃdvin} \textquotedblleft{}an adversary, a rival\textquotedblright{} (MW).
\end{cousinweb}

\begin{grammarlens}[title={What you've built}]
One more word for saying what happens if.
\end{grammarlens}

\subsection*{Your turn}
\begin{itemize}
\raggedright
  \item \emph{Say it:} \emph{pratidvandvī}
  \item \emph{Say it:} \emph{pratidvandvī}, once more
  \item \emph{Say it:} \emph{yadi pratidvandvī balavān, tarhi adhikaḥ abhyāsaḥ āvaśyakaḥ.} — If the rival is strong, more practice is needed.
  \item \emph{Recall:} say the Sanskrit for encourages, then the Sanskrit for competes, then say \emph{pratidvandvī} again
\end{itemize}

\begin{tcolorbox}[breakable,colback=teal!4,colframe=teal!35!black,title={Before you move on}]
What does \sk{प्रतिद्वन्द्वी} mean? (\textquotedblleft{}A rival\textquotedblright{}.) Say it once more.
\end{tcolorbox}
\section[\texorpdfstring{\sk{जेता} (jetā)}{जेता (jetā)}]{\sk{जेता} (jetā) — a winner; a victor}

\label{lesson:SA-C385-jeta}



\begin{quote}
Before the new one: say the Sanskrit for competes, then the Sanskrit for a rival.
\end{quote}

\subsection*{You'll want to know: \sk{जेता}}
\textbf{\sk{जेता}} — \emph{jetā} — \textquotedblleft{}a winner; a victor\textquotedblright{}.

\begin{quote}
\sk{यदि} \sk{भवान्} \sk{जेता} \sk{भवति}, \sk{तर्हि} \sk{पारितोषिकं} \sk{लभते।}
\end{quote}

\begin{quote}
\emph{yadi bhavān jetā bhavati, tarhi pāritoṣikaṃ labhate.}
\end{quote}

\begin{quote}
If you are the winner, you get the prize.
\end{quote}

A masculine noun.

\begin{cousinweb}
\textbf{\sk{जेता}} (\emph{jetā}) is built on the root √ji, \textquotedblleft{}to win\textquotedblright{}, with the agent ending \emph{-tṛ}; MW gives \emph{jetṛ} \textquotedblleft{}victorious, triumphant, a conqueror\textquotedblright{} (MW).
\end{cousinweb}

\begin{grammarlens}[title={What you've built}]
One more word for saying what happens if.
\end{grammarlens}

\subsection*{Your turn}
\begin{itemize}
\raggedright
  \item \emph{Say it:} \emph{jetā}
  \item \emph{Say it:} \emph{jetā}, once more
  \item \emph{Say it:} \emph{yadi bhavān jetā bhavati, tarhi pāritoṣikaṃ labhate.} — If you are the winner, you get the prize.
  \item \emph{Recall:} say the Sanskrit for competes, then the Sanskrit for a rival, then say \emph{jetā} again
\end{itemize}

\begin{tcolorbox}[breakable,colback=teal!4,colframe=teal!35!black,title={Before you move on}]
What does \sk{जेता} mean? (\textquotedblleft{}A winner; a victor\textquotedblright{}.) Say it once more.
\end{tcolorbox}
\section[\texorpdfstring{\sk{पुरस्करोति} (puraskaroti)}{पुरस्करोति (puraskaroti)}]{\sk{पुरस्करोति} (puraskaroti) — rewards (with a prize)}

\label{lesson:SA-C385-puraskaroti}



\begin{quote}
Before the new one: say the Sanskrit for a rival, then the Sanskrit for a winner.
\end{quote}

\subsection*{You'll want to know: \sk{पुरस्करोति}}
\textbf{\sk{पुरस्करोति}} — \emph{puraskaroti} — \textquotedblleft{}rewards (with a prize)\textquotedblright{}.

\begin{quote}
\sk{यदि} \sk{छात्रः} \sk{प्रथमः}, \sk{तर्हि} \sk{विद्यालयः} \sk{तं} \sk{पुरस्करोति।}
\end{quote}

\begin{quote}
\emph{yadi chātraḥ prathamaḥ, tarhi vidyālayaḥ taṃ puraskaroti.}
\end{quote}

\begin{quote}
If the student comes first, the school rewards him.
\end{quote}

Past (he …): \sk{सः} \sk{छात्रं} \sk{पुरस्कृतवान्}, \emph{saḥ chātraṃ puraskṛtavān}.

\begin{cousinweb}
\textbf{\sk{पुरस्करोति}} (\emph{puraskaroti}) is built from \emph{puras}, \textquotedblleft{}in front\textquotedblright{}, and the root √kṛ, \textquotedblleft{}to make\textquotedblright{}; MW gives \emph{puras-√kṛ} \textquotedblleft{}to place before\textquotedblright{} and \textquotedblleft{}to respect, to honour\textquotedblright{}, and \textquotedblleft{}rewards with a prize\textquotedblright{} is its modern sense (MW).
\end{cousinweb}

\begin{grammarlens}[title={What you've built}]
One more word for saying what happens if.
\end{grammarlens}

\subsection*{Your turn}
\begin{itemize}
\raggedright
  \item \emph{Say it:} \emph{puraskaroti}
  \item \emph{Say it:} \emph{puraskaroti}, once more
  \item \emph{Say it:} \emph{yadi chātraḥ prathamaḥ, tarhi vidyālayaḥ taṃ puraskaroti.} — If the student comes first, the school rewards him.
  \item \emph{Recall:} say the Sanskrit for a rival, then the Sanskrit for a winner, then say \emph{puraskaroti} again
\end{itemize}

\begin{tcolorbox}[breakable,colback=teal!4,colframe=teal!35!black,title={Before you move on}]
What does \sk{पुरस्करोति} mean? (\textquotedblleft{}Rewards (with a prize)\textquotedblright{}.) Say it once more.
\end{tcolorbox}
\section[\texorpdfstring{\sk{उत्कर्षः} (utkarṣaḥ)}{उत्कर्षः (utkarṣaḥ)}]{\sk{उत्कर्षः} (utkarṣaḥ) — excellence}

\label{lesson:SA-C385-utkarsah}



\begin{quote}
Before the new one: say the Sanskrit for a winner, then the Sanskrit for rewards (with a prize).
\end{quote}

\subsection*{You'll want to know: \sk{उत्कर्षः}}
\textbf{\sk{उत्कर्षः}} — \emph{utkarṣaḥ} — \textquotedblleft{}excellence\textquotedblright{}.

\begin{quote}
\sk{यदि} \sk{नित्यम्} \sk{अभ्यासं} \sk{करोति}, \sk{तर्हि} \sk{उत्कर्षः} \sk{निश्चितः।}
\end{quote}

\begin{quote}
\emph{yadi nityam abhyāsaṃ karoti, tarhi utkarṣaḥ niścitaḥ.}
\end{quote}

\begin{quote}
If one practises constantly, excellence is certain.
\end{quote}

A masculine noun.

\begin{cousinweb}
\textbf{\sk{उत्कर्षः}} (\emph{utkarṣaḥ}) is built on \emph{ut}, \textquotedblleft{}up\textquotedblright{}, and the root √kṛṣ, \textquotedblleft{}to draw\textquotedblright{}; MW gives \emph{ut-karṣa} \textquotedblleft{}excellence, eminence\textquotedblright{} (MW).
\end{cousinweb}

\begin{grammarlens}[title={What you've built}]
One more word for saying what happens if.
\end{grammarlens}

\subsection*{Your turn}
\begin{itemize}
\raggedright
  \item \emph{Say it:} \emph{utkarṣaḥ}
  \item \emph{Say it:} \emph{utkarṣaḥ}, once more
  \item \emph{Say it:} \emph{yadi nityam abhyāsaṃ karoti, tarhi utkarṣaḥ niścitaḥ.} — If one practises constantly, excellence is certain.
  \item \emph{Recall:} say the Sanskrit for a winner, then the Sanskrit for rewards (with a prize), then say \emph{utkarṣaḥ} again
\end{itemize}

\begin{tcolorbox}[breakable,colback=teal!4,colframe=teal!35!black,title={Before you move on}]
What does \sk{उत्कर्षः} mean? (\textquotedblleft{}Excellence\textquotedblright{}.) Say it once more.
\end{tcolorbox}
